\usepackage{ifthen}
%%%%%%%%%%%%%%%%%%%%%%%%%%%%%%%%%%%%%%%%%%%%%%%%%%%%%%%%%%%%%%%%%%%%%%%%%%%%%%%
%%  Modifiers and Fonts
%%%%%%%%%%%%%%%%%%%%%%%%%%%%%%%%%%%%%%%%%%%%%%%%%%%%%%%%%%%%%%%%%%%%%%%%%%%%%%%

\renewcommand{\k}[1]{\textit{#1}}
\newcommand{\mk}[1]{\mathit{#1}}

\renewcommand{\v}[1]{\texttt{#1}}
\newcommand{\mv}[1]{\text{\small\textsf{#1}}\xspace}
\newcommand{\mvs}[1]{\text{\scriptsize\textsf{#1}}\xspace}

\newcommand{\sref}[2]{\ref{#1}{\color{BrickRed}.}{\subref*{#1:#2}}}

%%%%%%%%%%%%%%%%%%%%%%%%%%%%%%%%%%%%%%%%%%%%%%%%%%%%%%%%%%%%%%%%%%%%%%%%%%%%%%%
%%  Symbols
%%%%%%%%%%%%%%%%%%%%%%%%%%%%%%%%%%%%%%%%%%%%%%%%%%%%%%%%%%%%%%%%%%%%%%%%%%%%%%%

% Tildes and crosses

\newcommand{\Y}{\ding{51}}
\newcommand{\N}{\ding{55}}

% Math

\newcommand{\card}[1]{\lvert#1\rvert}
\newcommand{\<}{\langle}
\renewcommand{\>}{\rangle}
\newcommand{\ldot}{\;.\;}
\newcommand{\lcom}{\,,\,}
\newcommand{\primeI}[1]{{#1}'}
\newcommand{\primeII}[1]{{#1}'\!\!\:'}
\newcommand{\into}{\mapsto}
\newcommand{\Nat}{\mathbb{N}}
\newcommand{\dotrightarrow}{\clipbox{0pt 0 {.675\width} 
0}{\ensuremath{\rightarrow}} \makebox[1.1pt]{} {\cdot} {\cdot} {\cdot} 
\makebox[1.1pt]{} \clipbox{{.55\width} 0 0pt -1pt}{\ensuremath{\rightarrow}}}
\newcommand{\longdotrightarrow}{\clipbox{0pt 0 {.55\width} 
0}{\ensuremath{\rightarrow}} \makebox[2pt]{} {\cdot} \makebox[1pt]{} {\cdot} 
\makebox[1pt]{} {\cdot} \makebox[2pt]{} \clipbox{{.45\width} 0 0pt 
0}{\ensuremath{\rightarrow}}}

% Sets

\newcommand{\C}{\subseteq}
\newcommand{\x}{\times}

% Logic

\newcommand{\true}{\mk{true}}
\newcommand{\false}{\mk{false}}
\renewcommand{\iff}{\Leftrightarrow}
\newcommand{\then}{\Rightarrow}

% Tools

\newcommand{\MTSA}{{\scshape mtsa}\xspace}
\newcommand{\SUP}{{\scshape supremica}\xspace}
\newcommand{\MBP}{{\scshape mbp}\xspace}
\newcommand{\PRP}{{\scshape prp}\xspace}
\newcommand{\MYND}{{\scshape mynd}\xspace}
\newcommand{\RATSY}{{\scshape ratsy}\xspace}
\newcommand{\ACACIA}{{\scshape acacia+}\xspace}
\newcommand{\SLUGS}{{\scshape slugs}\xspace}
\newcommand{\CIRCA}{{\scshape circa}\xspace}
\newcommand{\PARTY}{{\scshape party}\xspace}
\newcommand{\GPT}{{\scshape gpt}\xspace}
\newcommand{\DCS}{{\scshape otf-dcs}\xspace}
\newcommand{\DCSRA}{{\scshape otf-dcs-ra}\xspace}
%\newcommand{\DCSRAMON}{{\scshape mon-dcs-ra}\xspace}
\newcommand{\DCSBFS}{{\scshape otf-dcs-bfs}\xspace}
\newcommand{\DCSMA}{{\scshape otf-dcs-ma}\xspace}

% Supervisory Control

\renewcommand{\l}{\ell}
\renewcommand{\ll}{{\bar{\l}}}
\renewcommand{\d}{\delta}
\newcommand{\init}[1]{\bar{#1}} % #1 state name
\newcommand{\state}[2]{{#1}^{#2}}
\newcommand{\s}[1]{\state{e}{#1}}

\newcommand{\D}{\rightarrow}
\newcommand{\B}{\rightsquigarrow}

\newcommand{\A}{\mathcal{A}}
\newcommand{\E}{\mathcal{E}}
\newcommand{\G}{\mathcal{G}}
\newcommand{\M}{\mathcal{M}}
\newcommand{\R}[2]{{#1 \mathrel{R} #2}}
\newcommand{\W}{\mathcal{W}}

\renewcommand{\L}{\mathcal{L}}
\renewcommand{\P}{\mathcal{P}}
\renewcommand{\S}{\mathcal{S}}

\newcommand{\w}{\ensuremath{\omega}}

% Entornos de teoremas/definiciones (Jansen et al. 2019), usados en el Cap. 3
\theoremstyle{definition}
\newtheorem{definition}{Definición}[section]
\theoremstyle{plain}
\newtheorem{theorem}[definition]{Teorema}
\newtheorem{lemma}[definition]{Lema}
\newtheorem{proposition}[definition]{Proposición}

% Notación de branching bisimilarity (Jansen et al. 2019), usada en el Cap. 3
\newlength{\leftrightarrowwidth}
\settowidth{\leftrightarrowwidth}{$\leftrightarrow$}
\newcommand{\bisauxiliary}{\raisebox{.3ex}{\makebox[\leftrightarrowwidth]{$\underline{\makebox[0.7\leftrightarrowwidth]{$\leftrightarrow$}}$}}}
\newcommand{\bbis}{\mathrel{\bisauxiliary_b\!}}
\newcommand{\pijl}[1]{\mathrel{\text{$\xrightarrow{\smash[t]{#1}}$}}}
\newcommand{\lijp}[1]{\mathrel{\text{$\xleftarrow{\smash[t]{#1}}$}}}
% Slices y conjuntos derivados de las particiones Pi_s y Pi_t (Cap. 4)
\newcommand{\TB}{T_{B\pijl{}}}
\newcommand{\TBp}{T'_{B\pijl{}}}
\newcommand{\TBaB}{T_{B\pijl{a} B'}}
\newcommand{\TaB}{T_{\pijl{a} B'}}
\newcommand{\BT}{B_{\pijl{T}}}
% Ojo: \# esta redefinido mas abajo como \rapprox (notacion de DCS), asi que
% el numeral hay que tomarlo de la fuente de texto con \symbol{35}.
\newcommand{\hash}{\mathord{\text{\symbol{35}}}}
\newcommand{\nraB}{\hash aB'}
\newcommand{\Bottom}{\mathit{Bottom}}
\newcommand{\splittableBlocks}{\mathit{splittableBlocks}}

\newcommand{\notstep}[2]{\overset{#1}{\not\D}_{#2}}
\newcommand{\step}[2]{\overset{#1}{\D}_{#2}}
\newcommand{\mstep}[2]{\overset{#1}{\twoheadrightarrow}_{#2}}
\newcommand{\lstep}[2]{\overset{#1}{\longrightarrow}_{#2}}
\newcommand{\walk}[2]{\overset{#1}{\Rightarrow}_{#2}}
\newcommand{\hop}[2]{\overset{#1}{\B}_{#2}}
\newcommand{\runw}[2]{\,\overset{#1}{\dotrightarrow}_{#2}\,}
\newcommand{\run}[3]{
\,\overset{#1\ldots#2}{\longdotrightarrow}_{#3}\,}
\newcommand{\runlst}[3]{
\overset{#1\quad#2}{\raisebox{.1pt}{-}{\,\cdot\cdot\cdot}\!\D}_{#3}}
\newcommand{\enabled}[2]{\step{}{#1}\!\!(#2)}
\newcommand{\runs}[2]{\runw{}{#1}\!\!(#2)}

\newcommand{\trimlst}[1]{\trimbox{0 0 0 5pt}{\ensuremath{#1}}}

% DCS

\newcommand{\ma}[1]{\tilde{#1}}
\renewcommand{\#}{\text{\rapprox}}
\newcommand{\rapprox}{\protect\raisebox{.09em}{\:\!\protect\reflectbox{\protect\rotatebox[origin=c]{90}{\ensuremath{\approx}}}}}

\newcommand{\dist}[1]{\makebox[6pt][c]{\raisebox{0pt}[6pt][1.5pt]{\ensuremath{\scriptscriptstyle #1}}}}

\newcommand{\CCC}{\mathbb{C}}
\newcommand{\WCCC}{\mathbb{W}}

\newcommand{\open}{\mk{open}}
\newcommand{\recommendations}{\mk{RT}} %{\mathcal{R}} %\mk{recommendations}
\newcommand{\Goals}{\mk{Goals}}
\newcommand{\Errors}{\mk{Errors}}
\newcommand{\Witness}{\mk{Witnesses}}
\newcommand{\NONE}{\mk{None}}
\newcommand{\Unsettled}{\mk{Unsettled}}

%nuevos:  ------------------------------------------------------------
\newcommand{\pathBetween}[2]{(#1 \step{\l}{E} #2\runw{}{\structure} #1)}

\newcommand{\unexploredToBottom}[1]{\mk{#1_{\bot}}}
\newcommand{\unexploredToTop}[1]{\mk{#1_{\top}}}


%\newcommand{\unexploredToBottom}[1]{\mk{#1_{?\rightarrow \bot}}}
%\newcommand{\unexploredToTop}[1]{\mk{#1_{?\rightarrow \top}}}

\newcommand{\forcedTo}[3]{\mk{(#1\twoheadrightarrow #2)_{#3}}}

\newcommand{\WES}{\mk{W_{\unexploredToBottom{\structure}}}}
\newcommand{\WESS}{\mk{W_{\unexploredToBottom{\structure'}}}}
\newcommand{\LES}{\mk{L_{\unexploredToTop{\structure}}}}
\newcommand{\LESS}{\mk{L_{\unexploredToTop{\structure'}}}}
\newcommand{\WE}{\mk{W_E}}
\newcommand{\LE}{\mk{L_E}}

\newcommand{\SCC}{\mk{loops}} %strongly connected component

\newcommand{\Loops}[2]{AllNoneLoops(#1,#2)}

%\newcommand{\controllerReachsGoalsInOneStep}[1]{\mk{(\exists \l \ldot \trimlst{#1 \step{\l}{E} e''} \wedge e'' \in \Goals)    \wedge (\forall \l_u \in A_E^U \ldot \trimlst{#1 \step{\l_u}{E} e''} \then e'' \in \Goals)}}


\newcommand{\marked}[2]{\mk{\exists e_m \in #1 \ldot e_m \in M_{#2}}}
\newcommand{\Marked}{\mathit{Marked}}
\newtheorem{invariante}[definition]{Invariante}

\newcommand{\descendants}{\mk{descendants}}

\newcommand{\Targets}{Targets\newcommand{\Targets}{Targets}

}
%----------------------------------------------------------------------

\newcommand{\heuristic}{\mk{heuristic}}
\newcommand{\initial}{\ensuremath{\init{e}}}
\newcommand{\current}{\ensuremath{s}}
\newcommand{\child}{\ensuremath{s'}}
\newcommand{\structure}{\mk{ES}}
\newcommand{\toOpen}{\mk{toOpen}}

\newcommand{\ancestors}{\mk{ancestors}}
\newcommand{\predecesor}{\mk{predecesor}}
\newcommand{\statusu}{\mk{status}}
\newcommand{\visited}{\mk{visited}}
\newcommand{\pending}{\mk{pending}}
\newcommand{\unconfirmed}{\mk{unconfirmed}}

\newcommand{\target}{\mk{target}}
\newcommand{\novel}{\mk{novel}}
\newcommand{\distant}{\mk{distant}}
\newcommand{\strides}{\mk{strides}}

\newcommand{\MA}[2]{\mk{MA}_{#1}^{#2}}
\newcommand{\RA}[2]{\mk{RA}_{#1}^{#2}}

\newcommand{\g}{\ensuremath{g}}
\newcommand{\m}{m}%{\mathsf{m}}
\newcommand{\rr}{r}%{\mathsf{r}}
\newcommand{\pp}{p}%{\mathsf{p}}





\newboolean{showcomments}
\setboolean{showcomments}{true} % comment this line to 
%deactivate comments


% author macros ------------------------------------------------------------
\ifthenelse{\boolean{showcomments}}
{
	\newcommand{\ugh}[1]{\textcolor{red}{\uwave{#1}}}   
	% please rephrase
	\newcommand{\ins}[1]{\textcolor{blue}{\uline{#1}}}  % 
	%please insert
	\newcommand{\del}[1]{\textcolor{red}{\sout{#1}}}    % 
	%please delete
	\newcommand{\chg}[2]{
		\textcolor{red}{\sout{#1}}{$\rightarrow$}
		\textcolor{blue}{\uline{#2}}}               % please change
}{
	\newcommand{\ugh}[1]{#1}                            % please 
	%rephrase
	\newcommand{\ins}[1]{#1}                            % please insert
	\newcommand{\del}[1]{}                              % please delete
	\newcommand{\chg}[2]{#2}
}

\ifthenelse{\boolean{showcomments}}{
	\newcommand{\nbc}[3]{
		{\colorbox{#3}{\bfseries\sffamily\scriptsize\textcolor{white}{#1}}}
		{\textcolor{#3}{\sf\small$\langle$\textit{#2}$\rangle$}}}
}{
	\newcommand{\nbc}[3]{}
	\newcommand{\version}{}
}

\newcommand\su[1]{\nbc{SU}{#1}{red}} 
\newcommand\cl[1]{\nbc{CL}{#1}{pink}} 
\newcommand\hg[1]{\nbc{HG}{#1}{olive}}
\newcommand\vb[1]{\nbc{VB}{#1}{teal}}

\newcommand\resolved[1]{}
% add more author macros here

\newcommand\tocado[1]{\textcolor{red}{#1}}
\newcommand\comment[1]{\textcolor{cyan}{#1}}

\newcommand{\multlinecomment}[1]{\directlua{-- #1}}

\newcommand{\Review}[1]{#1}

%%%%%%%%%%%%%%%%%%%%%%%%%%%%%%%%%%%%%%%%%%%%%%%%%%%%%%%%%%%%%%%%%%%%%%%%%%%%%%%
%%  Notacion de complejidad y llaves de costo en los algoritmos (Cap. 4)
%%  Tomados de las fuentes de Jansen et al. (arXiv:1909.10824).
%%%%%%%%%%%%%%%%%%%%%%%%%%%%%%%%%%%%%%%%%%%%%%%%%%%%%%%%%%%%%%%%%%%%%%%%%%%%%%%

\newcommand{\bigo}[1]{{O\!\left(#1\right)}}
\newcommand{\bigO}{O}                                   % O sin argumento
\newcommand{\size}[1]{\left\lvert #1 \right\rvert}

% \rightbrace{<lineas>}{<texto>}: llave gris al margen derecho de un bloque de
% pseudocodigo, para anotar su costo.
\newlength{\halflineskip}
\newcommand{\rightbrace}[3][0]{\setlength{\halflineskip}{0.5\baselineskip}\hspace{0pt plus 1filll}\makebox[0pt][l]{\smash{\tikz[baseline=#1\baselineskip]{\draw[thick,color=gray,line cap=round,decorate,decoration={brace,amplitude=2.3pt}] (0pt,0.7\baselineskip) ++(0.1em,0) -- ++(0pt,-#2\baselineskip+1pt) (5pt,\halflineskip-#2\halflineskip) node[color=gray,right,inner sep=0.1em,anchor=base west]{\small #3};}}}}

% Comentarios en gris dentro de los algoritmos
\renewcommand{\algorithmiccomment}[1]{\hfill{\color{gray}// #1}}

% Rotulo del float de algoritmos ("Algorithm 1" -> "Algoritmo 1")
\makeatletter
\renewcommand{\ALG@name}{Algoritmo}
\makeatother
\renewcommand{\listalgorithmname}{Índice de algoritmos}

% Palabras clave del pseudocodigo en castellano. Las compuestas
% (\algorithmicendif, \algorithmicendfor, \algorithmicendwhile, ...) se derivan
% de \algorithmicend y de la palabra correspondiente, asi que salen solas.
\renewcommand{\algorithmicrequire}{\textbf{Requiere:}}
\renewcommand{\algorithmicensure}{\textbf{Asegura:}}
\renewcommand{\algorithmicend}{\textbf{fin}}
\renewcommand{\algorithmicif}{\textbf{si}}
\renewcommand{\algorithmicthen}{\textbf{entonces}}
\renewcommand{\algorithmicelse}{\textbf{si no}}
\renewcommand{\algorithmicelsif}{\textbf{si no, si}}
\renewcommand{\algorithmicfor}{\textbf{para}}
\renewcommand{\algorithmicforall}{\textbf{para todo}}
\renewcommand{\algorithmicdo}{\textbf{hacer}}
\renewcommand{\algorithmicwhile}{\textbf{mientras}}
\renewcommand{\algorithmicloop}{\textbf{repetir}}
\renewcommand{\algorithmicrepeat}{\textbf{repetir}}
\renewcommand{\algorithmicuntil}{\textbf{hasta que}}
\renewcommand{\algorithmicprint}{\textbf{imprimir}}
\renewcommand{\algorithmicreturn}{\textbf{retornar}}
\renewcommand{\algorithmicand}{\textbf{y}}
\renewcommand{\algorithmicor}{\textbf{o}}
\renewcommand{\algorithmicnot}{\textbf{no}}
\renewcommand{\algorithmicto}{\textbf{hasta}}
\renewcommand{\algorithmicinputs}{\textbf{entradas}}
\renewcommand{\algorithmicoutputs}{\textbf{salidas}}
\renewcommand{\algorithmicglobals}{\textbf{globales}}
\renewcommand{\algorithmicbody}{\textbf{hacer}}
\renewcommand{\algorithmictrue}{\textbf{verdadero}}
\renewcommand{\algorithmicfalse}{\textbf{falso}}

%%%%%%%%%%%%%%%%%%%%%%%%%%%%%%%%%%%%%%%%%%%%%%%%%%%%%%%%%%%%%%%%%%%%%%%%%%%%%%%
%%  Estilos tikz para las figuras de bloques y bunches (Cap. 4)
%%  Tomados de las fuentes de Jansen et al. (arXiv:1909.10824).
%%  Van agrupados en el estilo `bunchpic' para que la redefinicion de `state'
%%  quede local a cada tikzpicture y no afecte a los automatas de los otros
%%  capitulos (que usan el `state' de la libreria automata).
%%%%%%%%%%%%%%%%%%%%%%%%%%%%%%%%%%%%%%%%%%%%%%%%%%%%%%%%%%%%%%%%%%%%%%%%%%%%%%%

\newlength{\statetextwidth}\settowidth{\statetextwidth}{$s_0$}
\newlength{\statetextheight}\settoheight{\statetextheight}{$s_0$}
\newlength{\statetextdepth}\settodepth{\statetextdepth}{$s_0$}

\tikzset{bunchpic/.style={
	bunch/.style={very thin,draw,fill=gray!5},
	block/.style={thick,draw,fill=gray!15,rounded corners=4.95mm},
	state/.style={circle,draw,semithick,fill=white,inner sep=1.35pt,
		text height=\statetextheight,text depth=\statetextdepth},
	% Estados rayados: verticales para los destinados a R, horizontales para
	% los destinados a U. Se usan rayas (y no solo color) para que las figuras
	% se sigan entendiendo impresas en blanco y negro.
	red state/.style={state,fill=red!10,path picture={\draw[red!40,very thick]
		(-0.25\statetextwidth,-1cm) -- (-0.25\statetextwidth,2cm)
		(0cm,-1cm) -- (0cm,2cm)
		(0.25\statetextwidth,-1cm) -- (0.25\statetextwidth,2cm)
		(0.5\statetextwidth,-1cm) -- (0.5\statetextwidth,2cm)
		(0.75\statetextwidth,-1cm) -- (0.75\statetextwidth,2cm)
		(\statetextwidth,-1cm) -- (\statetextwidth,2cm)
		(1.25\statetextwidth,-1cm) -- (1.25\statetextwidth,2cm);}},
	blue state/.style={state,fill=blue!10,path picture={\draw[blue!40,very thick]
		(-1cm,0.5\statetextheight-0.5\statetextdepth+0.75\statetextwidth) -- (2cm,0.5\statetextheight-0.5\statetextdepth+0.75\statetextwidth)
		(-1cm,0.5\statetextheight-0.5\statetextdepth+0.5\statetextwidth) -- (2cm,0.5\statetextheight-0.5\statetextdepth+0.5\statetextwidth)
		(-1cm,0.5\statetextheight-0.5\statetextdepth+0.25\statetextwidth) -- (2cm,0.5\statetextheight-0.5\statetextdepth+0.25\statetextwidth)
		(-1cm,0.5\statetextheight-0.5\statetextdepth) -- (2cm,0.5\statetextheight-0.5\statetextdepth)
		(-1cm,0.5\statetextheight-0.5\statetextdepth-0.25\statetextwidth) -- (2cm,0.5\statetextheight-0.5\statetextdepth-0.25\statetextwidth)
		(-1cm,0.5\statetextheight-0.5\statetextdepth-0.5\statetextwidth) -- (2cm,0.5\statetextheight-0.5\statetextdepth-0.5\statetextwidth)
		(-1cm,0.5\statetextheight-0.5\statetextdepth-0.75\statetextwidth) -- (2cm,0.5\statetextheight-0.5\statetextdepth-0.75\statetextwidth);}}
}}

%%  Marcas en linea para referirse a los estados rayados desde el texto.
\newcommand{\redstatetext}{\smash{\tikz[bunchpic,baseline=-0.37\statetextwidth]{%
	\node[red state,thin,inner sep=0pt]{\makebox[0.5\statetextwidth]{}};}}}
\newcommand{\bluestatetext}{\smash{\tikz[bunchpic,baseline=-0.37\statetextwidth]{%
	\node[blue state,thin,inner sep=0pt]{\makebox[0.5\statetextwidth]{}};}}}